\section{Transfer of the neutral elliptic construction}
\label{ch:ell:section}

The charged construction uses the neutral system with the total tensors
$(g,K)$. We verify its prepared class and identify the equations,
normalizations and estimates needed from
Theorem~\ref{n3:existence:deformation}, Lemmas~\ref{n3:apriori:collars}
and~\ref{n3:limit:end}, with the parameter choices in
Proposition~\ref{n3:deformation:floor}. The transfer proposition below is a
consequence of those proved results; it introduces no new solvability
hypothesis.

\subsection{The class to which the input applies}

Fix one strict rest datum supplied by Proposition~\ref{ch:rest:reduction},
with original obstacle $S_0$, energy $E_*>0$, and
$A_*=a_g(S_0)$.  We record the complete preparation relevant to the
neutral construction, so that its hypotheses are not replaced by a solvability
assumption about arbitrary charged data.

\begin{definition}[Prepared neutral class]\label{ch:ell:prepared-neutral}
The underlying reduced datum is smooth, orientable, connected, complete
up to its nonempty compact smooth boundary $S_0$, and has exactly one
asymptotically flat end.  Its boundary satisfies
$H_{S_0}+\tr_{S_0}K<0$, with normal toward that end.  The total
constraint densities satisfy
\[
 \mu,|J|_g\in L^1(dV_g),\qquad
 M_N:=8\pi(\mu-|J|_g)>0,\qquad
 M_N\ge c r_y^{-3-\delta}\quad\hbox{far out},
 \qquad 0<\delta<1,
\]
for a fixed $c>0$.  The tensor $K$ is compactly supported.  In a rest
chart $y$, $g-\delta_{\rm eucl}=O_j(r_y^{-1})$ for every fixed $j$,
and hence $\Ric_g=O(r_y^{-3})$; the ADM energy is finite and the
momentum is zero.  Our reduced data also have
$R_g=O_j(r_y^{-3-\delta})$ for every fixed $j$.

The prepared exterior $X$ is the closure of the component toward
infinity remaining after the full black region and the subsequent full
white region in Proposition~\ref{ch:prep:prepared} are removed.  More
precisely, choose $c_b<0$ with
$c_b>\max_{S_0}(H+\tr_{S_0}K)$ as a Hausdorff right-continuity point
of the full maximal family with $H+\tr_\Sigma K\le c_b$ and the required
inner collar of $S_0$.  With that black region fixed, choose $c_w<0$
as a right-continuity point for the full white family with
$H-\tr_\Sigma K\le c_w$, no required inner boundary, and the black faces
and a distant sphere as outer barriers.  If the selected region is
empty, right continuity includes emptiness immediately to its right.
The adjacent frontiers give a smooth compact boundary
$B=B_b\sqcup B_w\ne\varnothing$, with
\begin{equation}\label{ch:ell:face-signs}
 H+P_B=c_b\quad\hbox{on }B_b,\qquad
 H-P_B=c_w\quad\hbox{on }B_w,
 \qquad P_B=\tr_BK,
\end{equation}
where $\nu$ points into $X$ and $H=\Div_B\nu$.
The faces have disjoint smooth outward leaf collars, with leaf
coordinate $s\ge0$, $|ds|>0$, and the expansion of the corresponding
color strictly greater than its threshold for $s>0$.
Every cut enclosing $B$ has $g$-area at least $A_*$.

The fixed offset and weights satisfy
\begin{gather*}
 C\in C_c^\infty(\overline X),\qquad C_w\le C\le C_b,
 \qquad C_b\ge0,\quad C_w\le0,\\
 C=C_b-k_Bs\ \hbox{near a black face},\qquad
 C=C_w+k_Bs\ \hbox{near a white face},\qquad k_B>0,\\
 b_-=c_b-C_b<0,\qquad b_+=-c_w-C_w>0.
\end{gather*}
Here $k_B$ may depend on the face.  Either color may be absent; its
constant $C_b$ or $C_w$ then remains only a bound, and its face
conditions are omitted.  The function $\rho$ is smooth and positive,
equals a positive constant times $r_y^{-4}$ sufficiently far out,
and has symbol derivative bounds there.  For every
$h\in[b_-,b_+]$, the exact neutral smallness condition is
\begin{equation}\label{ch:ell:neutral-smallness}
 |(h+C)\tr K|+|\nabla C|+\rho\le\tfrac12 M_N
 \quad\hbox{throughout }X.
\end{equation}
All thresholds, full regions, collars, offsets, and these weights are
fixed before any deformation parameter is chosen.  This is the
preparation in Proposition~\ref{n3:domains:prepared-domain} on the reduced
class of Proposition~\ref{n3:reduction:rest-reduction}.
\end{definition}

Here ``reduced class'' means the properties in
Proposition~\ref{n3:reduction:rest-reduction}\textup{(i)--(ii)}, which
explicitly singles out these properties as the class for its subsequent
argument. The preparation uses these properties of the datum and the
small offset choices of Section~\ref{n3:domains:section} and
Proposition~\ref{n3:domains:prepared-domain}.

\begin{lemma}[Verification for the charged preparation]
\label{ch:ell:hypothesis-check}
The outputs of Propositions~\ref{ch:rest:reduction} and
\ref{ch:prep:prepared} belong to Definition~\ref{ch:ell:prepared-neutral}.
\end{lemma}
\begin{proof}
Proposition~\ref{ch:rest:reduction} supplies the reduced geometry,
integrability and end estimates; Proposition~\ref{ch:prep:prepared}
supplies the full-region preparation. Equation~\eqref{ch:pre:neutral-margin}
gives $M_N\ge M_0$, so the strict physical margin and its end lower
bound imply their neutral counterparts. The stronger bound
\eqref{ch:prep:offset-smallness} implies \eqref{ch:ell:neutral-smallness}.
These are the actual total constraint densities of $(g,K)$, as
required by the neutral construction.
\end{proof}

\subsection{Profiles, quadratic form, and boundary value problem}

All derivatives and contractions below use $g$; vectors and covectors
are identified by $g$.  Extend $r_y$ smoothly over $X$, bounded below
by one, and let $X_R$ be a large spherical truncation, with outer face
$S_R$.  First fix a sufficiently small $0<\epsilon<1$, then a large
integer $n\ge4$ with $n>2/\epsilon$, and then a sufficiently large
$R$.  These are the standing parameter conventions of
Section~\ref{n3:deformation:section}.  Choose a smooth nonincreasing
$\vartheta:\R\to[0,1]$, equal to one on $(-\infty,0]$ and zero
on $[1,\infty)$, and set
\begin{equation}\label{ch:ell:profiles}
 \ell=e^{-\epsilon n},\qquad \tau_0=\ell^{3/2},\qquad
 p(t)=n\vartheta(nt),\qquad
 l(t)=\exp\!\left(\int_0^t p(s)\,ds\right),\qquad
 L(t)=e^{2t}l(t).
\end{equation}
The symbol $\tau_0$ is unrelated to $\tau=\tr K$.
In particular $l(t)=e^{nt}$ for $t\le0$, $l$ is constant for
$t\ge1/n$, and $\ell\le l\le e$ when $t\ge-\epsilon$.
Define
\begin{equation}\label{ch:ell:variables}
 \begin{gathered}
 h=\tau_0f,\quad \sigma=|\nabla f|,\quad
 D=1+l^2\sigma^2,\quad d=D^{-1},\quad
 w=\frac{l\nabla f}{\sqrt D},\quad a=|w|,\quad v=a^2=1-d,\\
 \chi=\frac{4d}{4+pv},\qquad
 A_d=I-w\otimes w,\qquad
 A_\chi=I-\frac{p+4}{4+pv}w\otimes w,\\
 Z=t+\tfrac18\log D,\qquad u=L\sqrt D,\qquad
 H^f=\frac l{\sqrt D}\Hess f,\qquad S_f=K+H^f,\\
 F=\tr_{A_\chi}S_f,\qquad
 V=uA_\chi\bigl(4\nabla Z+K(w,\cdot)\bigr).
 \end{gathered}
\end{equation}
Matrix expressions are in a $g$-orthonormal frame, and
$\tr_A S=A^{ij}S_{ij}$.  The unknowns are $(f,Z)$:
for each fixed $\nabla f$, the displayed expression for $Z$ is
strictly increasing and onto as a function of $t$, since
$\partial_tZ=1+pv/4>0$.  Thus it determines $t$ smoothly without
assuming a floor.  In particular $0<\chi\le d\le1$.

For $\sigma>0$, put $e=\nabla f/\sigma$ and decompose
\[
 T=S_f|_{e^\perp\times e^\perp},\quad
 M=S_f(e,\cdot)|_{e^\perp},\quad k=S_f(e,e),\quad
 y_t=(\nabla t)_\perp,\quad x_t=\partial_et.
\]
These $T,M,k$ are tensor blocks; $M$ is not a manifold or an ADM
mass.  We have $\tr T=F-\chi k$, and $T^{\rm tf}$ means the
trace-free part in the two-dimensional plane $e^\perp$.  Set
\begin{equation}\label{ch:ell:quadratic}
 \begin{split}
 \cT={}&\tfrac12|T^{\rm tf}|^2
       +|M+(p+1)ay_t|^2+[4(p+1)-v]|y_t|^2\\
 &+4(p+1)d\left(x_t+\frac{\chi ak}{4d}\right)^2
       +\tfrac34\left(F-\frac{\chi k}{3}\right)^2
       +\frac{4d(9-d)}{3(4+pv)^2}k^2.
 \end{split}
\end{equation}
This is the exact square completion in Section~\ref{n3:deformation:section}.
It extends smoothly across $\sigma=0$; its direction-independent
value there is
\[
 \cT=\tfrac12\bigl(|S_f|^2+(\tr S_f)^2\bigr)
                  +4(p+1)|\nabla t|^2.
\]
The invariant expression and smooth extension are proved in
Lemma~\ref{ch:cmp:scalar}, and polynomial coercivity is proved in
Lemma~\ref{ch:cmp:coercivity}.

Fix $3/4<\gamma<1$ and positive smooth weights with
$\rho_0\asymp r_y^{-4}$, $\rho_1\asymp r_y^{-2\gamma}$ and
symbol derivative bounds on the end.  In particular
$\rho_0,\rho_0^2,\rho_1^2$ are integrable; integrability of
$\rho_1$ itself is not required.  For a smooth
$m_c:\R\to[0,1]$, equal to one on $(-\infty,0]$ and zero on
$[1,\infty)$, put
\[
 m_0(t)=m_c(n(t+\epsilon)),\qquad
 \cP=C_n(\rho_0+v\rho_1).
\]
Choose $N_B\in C^\infty(B)$ with $N_B>1+\sup_B|H|$.
Use the choices furnished by Lemma~\ref{n3:deformation:floor-errors} and Proposition~\ref{n3:deformation:floor}:
first a sufficiently small $0<\delta_0<1$, independent of $n,R$,
then a sufficiently large polynomial penalty
$C_n=C_*(1+n)^{N_*}$.  The fixed constants may depend on the
prepared data and auxiliary profiles.  All auxiliary collar cutoffs,
face extensions, and continuation choices are those furnished in
that construction.  Its physical endpoint is exactly
\begin{equation}\label{ch:ell:system}
 \begin{cases}
 F=h+C &\text{in }X_R,\\
 \Div V=u\bigl(\delta_0\cT+\rho+
           \delta_0\tau_0a\sigma-m_0(t)\cP\bigr)
       &\text{in }X_R,\\
 h=b_-\ \text{on }B_b,\qquad h=b_+\ \text{on }B_w,\\
 V_\nu/u=-H+w_\nu(F-P_B)-N_Bd &\text{on }B,\\
 f=Z=0 &\text{on }S_R.
 \end{cases}
\end{equation}
Here $V_\nu=g(V,\nu)$ and $w_\nu=g(w,\nu)$; the inner normal
is the negative of the integration-outward normal of $X_R$.
At the assigned face heights,
$F=P_B+H$ on black faces and $F=P_B-H$ on white faces.
On the admissible penalty band one has
\begin{equation}\label{ch:ell:floor-scale}
 -\epsilon\le t\le-\epsilon+1/n<0,
 \qquad \ell\le l\le e\ell,
 \qquad c\ell\le L\le C\ell.
\end{equation}
The last constants can be uniform for $0<\epsilon<1$.

\subsection{The transferred existence and estimates}

\begin{proposition}[Neutral elliptic construction for the prepared class]\label{ch:ell:neutral-input}
For the prepared class in Definition~\ref{ch:ell:prepared-neutral},
with the profiles, parameters, penalty choices, and physical endpoint
in Equations~\eqref{ch:ell:profiles}--\eqref{ch:ell:system}, the following
statements hold.
\begin{enumerate}
\item \textup{Existence with the floor bound.}
With $\delta_0,C_n$ chosen as in Proposition~\ref{n3:deformation:floor},
for every sufficiently small fixed $\epsilon>0$,
Theorem~\ref{n3:existence:deformation} supplies $n_0$, enlarged to include
$n\ge4$ and $n>2/\epsilon$, and a sufficient radius
$R_{\min}(n,\epsilon)\to\infty$ as $n\to\infty$, such that
Equation~\eqref{ch:ell:system} has a smooth solution on every $X_R$
with $n\ge n_0$ and $R\ge R_{\min}(n,\epsilon)$.
It satisfies $t>-\epsilon$ on $\overline{X_R}$.
The radius includes the geometric supports, all a priori and floor
thresholds, and the outer-boundary floor restriction in the proof of
that theorem.  No polynomial bound for this radius is asserted.

\item \textup{Physical-endpoint height and slope bounds.}
For smooth finite-truncation solutions with $t\ge-\epsilon$, for
these sufficiently large $n,R$, Lemma~\ref{n3:apriori:collars} at the
physical endpoint gives fixed $c,C_*>0$ such that
\begin{equation}\label{ch:ell:slopes}
 \frac c{\tau_0}\le\partial_\nu f\le\frac{C_*}{\tau_0}
       \quad(B_b),\qquad
 \frac c{\tau_0}\le-\partial_\nu f\le\frac{C_*}{\tau_0}
       \quad(B_w).
\end{equation}
These constants are independent of $n,R$ after fixing the prepared
data and $\epsilon$.  Also $b_-\le h\le b_+$ on every fixed
compact region needed for the scalar-curvature argument, in particular
on the supports of $K,C$.  A condition on an absent face color is
vacuous.  No global height bound is being added to this statement.

\item \textup{Exhaustion and end normalization.}
At fixed $\epsilon,n$, the solutions have local smooth estimates of
every fixed order, uniformly as $R\to\infty$, including up to the
fixed inner boundary.  Every sequence of allowed radii tending to
infinity has a subsequence converging smoothly on compact subsets of
$\overline X$ to a solution of Equation~\eqref{ch:ell:system} with
the outer boundary condition omitted and $t\ge-\epsilon$.
The preceding compact height and face estimates pass to this limit.
For sufficiently large $n\ge\max\{4,\lceil2/\epsilon\rceil\}$,
choose an exhausting sequence and the limiting solution supplied by
Lemma~\ref{n3:limit:end}.  For some $c_n>0$ and constants $C_j(n)$,
this solution also satisfies
\begin{equation}\label{ch:ell:end-estimates}
 |\nabla^jf|\le C_j(n)e^{-c_nr_y}\quad\text{for each fixed }j,
 \qquad t,Z=O_2(r_y^{-1}).
\end{equation}
Writing $\bar g=g+l^2df^2$, both $t-Z$ and $\bar g-g$, together
with their coordinate derivatives of every fixed order, decay
exponentially.  Moreover
\begin{equation}\label{ch:ell:end-flux}
 V=4\nabla_gt+O(r_y^{-3}),\qquad
 \lim_{s\to\infty}\int_{r_y=s}g(V,\nu_s)\,dA_g
 \quad\text{exists and is finite},
\end{equation}
where $\nu_s$ points toward infinity.
\end{enumerate}
\end{proposition}

\begin{proof}
We first identify the input class. The reduced properties in
Definition~\ref{ch:ell:prepared-neutral} are the properties
\textup{(i)--(ii)} of Proposition~\ref{n3:reduction:rest-reduction};
that proposition explicitly designates them as the class for the
subsequent construction. In particular, provenance as a particular
member of its approximation sequence is not an additional hypothesis.
The full black and white regions, their right-continuity conventions
including emptiness, the adjacent complete boundary components, and the
disjoint outward collars are exactly those of
Proposition~\ref{n3:domains:prepared-domain}. Equation~\eqref{ch:ell:neutral-smallness}
is its offset inequality with $M_N=8\pi(\mu-|J|)$.
Thus its preparation, including the possibility of either missing face
color, applies to every datum in the stated class.

Next identify the system. The neutral construction uses the capillarity
parameter denoted here by $\tau_0=\ell^{3/2}$, to distinguish it from
$\operatorname{tr}K$. Its functions $l,L$, logarithmic profile derivative
$p=l'/l$, graph quantities $D,d,w,a,v$, matrices $A_d,A_\chi$, and
unknowns $(f,Z)$ are exactly Equations~\eqref{ch:ell:profiles}--\eqref{ch:ell:variables}.
The notation $p_L=L'/L$ in the neutral construction therefore means
$p_L=p+2$. Its tensor $S$ and axial component $q$ are denoted here by
$S_f$ and $k$, respectively. Its quadratic form is
Equation~\eqref{ch:ell:quadratic},
with no change of sign in the momentum constraint. The physical endpoint
of Proposition~\ref{n3:deformation:homotopy} is its first continuation stage
at $s=0$, with drift coefficient $\lambda=1$: its trace equation is
$F=\tau_0f+C$, and
its divergence source is
$\delta_0\mathcal T+\rho+\delta_0\tau_0a\sigma-m_0(t)\mathcal P$.
These are the two interior equations of Equation~\eqref{ch:ell:system}.
The inner normal in both formulations points toward the end. The
componentwise assigned heights and the identities
$F=P_B+H$ on black faces and $F=P_B-H$ on white faces give exactly the
same inner flux condition; the outer values are $f=Z=0$ in both systems.
All auxiliary continuation data are chosen by that construction after
the fixed preparation, as stipulated above.

Lemma~\ref{n3:deformation:floor-errors} and
Proposition~\ref{n3:deformation:floor} choose a fixed $\delta_0>0$ and a
polynomial $C_n$. Theorem~\ref{n3:existence:deformation} then supplies item
\textup{1}, including its sufficiently large truncation radius and the
strict floor; enlarging the lower bounds on $n$ and $R$ imposes the
additional displayed numerical restrictions. At the physical endpoint,
Lemma~\ref{n3:apriori:collars} gives item \textup{2}. Its slope constants
are independent of $n,R$, its height interval is asserted on the fixed
compact regions used in the curvature argument, and its face conditions
are componentwise. No boundary connectedness is required.

For fixed $\epsilon,n$, Proposition~\ref{n3:existence:classical-bounds}
and the exhaustion following Theorem~\ref{n3:existence:deformation} give
the local smooth estimates and subsequential limit in item \textup{3},
including every fixed inner face. These estimates allow arbitrary
fixed-$n$ constants and do not change the polynomial penalty already
chosen. Lemma~\ref{n3:limit:end} supplies the stated normalized end solution:
the exponential decay of $f$ and the graph errors, the $O_2(r_y^{-1})$
bounds for $t,Z$, and the finite flux with
$V=4\nabla_gt+O(r_y^{-3})$. These are exactly
Equations~\eqref{ch:ell:end-estimates}--\eqref{ch:ell:end-flux}.
All three conclusions follow from the proved neutral construction.
\end{proof}

The constants in the local regularity and end estimates may depend
arbitrarily on the fixed $n$ (and on the prepared datum and
$\epsilon$); neither they nor $c_n^{-1}$ are claimed polynomial.
The penalty $C_n$, by contrast, is polynomial in $n$, as are the
explicit algebraic losses proved in Lemma~\ref{ch:cmp:coercivity}.
No uniqueness or convergence as $n\to\infty$ is part of
Proposition~\ref{ch:ell:neutral-input}.

The neutral results used in the transfer have no Penrose equality premise.
Proposition~\ref{ch:ell:neutral-input} uses their original physical
endpoint, corresponding to the first continuation stage with its
drift coefficient equal to one.  It introduces neither a charged
source term nor a solver for electromagnetic-subtracted constraints.
Lemma~\ref{ch:ell:hypothesis-check} verifies its hypotheses for each
fixed charged preparation.  Henceforth we use its global limiting
solution, always taking $R\to\infty$ first at fixed $\epsilon,n$.
The later numerical argument takes $n\to\infty$, then
$\epsilon\downarrow0$, and only then removes the strict rest-data
approximation.
